% ============================================================================
%  MSc Robotics and Computation — summer project report
%
%  Structure scaffold only. The prose is deliberately NOT written: the module
%  guidance (COMP0247 §5.5) is explicit that AI-drafted reports read shallow and
%  repetitive and mark poorly, and that any AI assistance must be credited. What
%  is here is the skeleton, the page budget, and notes on what each section has
%  to establish — the argument and the words should be yours.
%
%  One-column report class rather than a two-column article format, because the
%  results tables are wide (4-6 policy arms x TSR/collision/CI/ETS), there is
%  real display math, and the 30-page cap applies either way.
%
%  Limits (COMP0247 §5): 30 pages main body, 60 pages including bibliography and
%  appendices. Readable text size, reasonable margins, numbered pages.
% ============================================================================
\documentclass[a4paper,11pt]{report}

\usepackage[a4paper, margin=3cm, bottom=2.5cm]{geometry}
\usepackage{setspace}
\usepackage[T1]{fontenc}
\usepackage[utf8]{inputenc}

% --- maths -----------------------------------------------------------------
\usepackage{amsmath,amssymb,amsfonts,bm}

% --- figures and tables ----------------------------------------------------
\usepackage{graphicx}
\usepackage{booktabs}      % \toprule/\midrule/\bottomrule — never use \hline
\usepackage{siunitx}       % align numbers on the decimal point in results tables
\usepackage{subcaption}
\usepackage{multirow}
\usepackage{threeparttable}  % table notes tied to the table, for tab:external_validation

% Draft-mode placeholders for assets not yet collected. Set \figuretodofalse before submission:
% every placeholder box vanishes and any figure still missing becomes a visible ?? instead of a
% silently absent one.
\newif\iffiguretodo
\figuretodotrue
\graphicspath{{figures/}}

% --- references ------------------------------------------------------------
\usepackage[numbers,sort&compress]{natbib}   % \citet{x} -> "Author et al. [1]"
\usepackage[hidelinks]{hyperref}
\usepackage[capitalise,noabbrev]{cleveref}   % \cref{fig:x} -> "Figure 3.2"

% siunitx setup for percentage-with-CI columns
\sisetup{detect-weight=true, detect-inline-weight=math}

\newcommand{\pizero}{$\pi_{0.5}$}
\newcommand{\car}{\textsc{car}}
\newcommand{\tsr}{\textsc{tsr}}
\newcommand{\ets}{\textsc{ets}}

% ============================================================================
\title{{\vspace{-14em}
\includegraphics[scale=0.4]{ucl_logo.png}}\\
  {\Huge Project Title}\\
  {\large Optional Subtitle}\\}
\date{Submission date: Day Month Year}
\author{Candidate number: \#\#\#\#\#   % candidate number, NOT your name
  \thanks{{\bf Disclaimer:} This report is submitted as part requirement for the
    MSc Robotics and Computation at UCL. It is substantially the result of my own
    work except where explicitly indicated in the text.
    % KEEP ONLY ONE of the two disclaimers below.
    The report may be freely copied and distributed provided the source is
    explicitly acknowledged.}
  \\ \\ MSc Robotics and Computation \\ \\ Supervisor: NAME}

\begin{document}
\onehalfspacing
\maketitle

\begin{abstract}
% ~200 words. Write this LAST. Motivation, method, headline numbers.
% Your headline is a shield-free safety result plus the CBF-activation-rate drop.
\end{abstract}

\tableofcontents
\setcounter{page}{1}

% ============================================================================
\chapter{Introduction}\label{ch:intro}
% BUDGET ~3 pages.
% Establish, in order:
%  - VLAs can follow instructions but have no notion of collision safety.
%  - Why an external safety filter is the standard answer, and why that is
%    unsatisfying: it is always needed, and it costs task success.
%  - The question this project actually asks: can the safety be moved INTO the
%    policy so the filter becomes redundant?
%  - Scope: SafeLIBERO, pi_0.5, simulation only. Say what is NOT in scope.
%  - Contributions, as a short list. Be honest about which are negative results.

\section{Motivation}
\section{Problem statement and scope}
\section{Contributions}

% ============================================================================
\chapter{Background and Related Work}\label{ch:background}
% BUDGET ~5 pages. Concision is explicitly marked for — do not pad this.
% Must establish that your approach is the logical next step, not a survey.

\section{Vision-language-action models}
% pi_0.5, flow-matching action head, action chunking. Only the detail a reader
% needs to follow your method.

\section{Control barrier functions for manipulation}
% h(x) >= 0, the QP projection, why it composes with any policy.

\section{Safety filters for VLAs}
% AEGIS is the direct comparison — state exactly what they do and where they
% stop (their documented distribution-shift limitation is your opening).

\section{Imitation and RL for safety internalisation}
% DAgger, SafeDAgger, Flow-GRPO, piRL. Position: piRL is task RL for flow VLAs;
% this project is safety RL / safety distillation.

% ============================================================================
\chapter{Method}\label{ch:method}
% BUDGET ~8 pages. This is where the marks for "difficulty and achievement" are
% won. Use figures — a system diagram and a barrier-geometry figure.

\section{The SafeLIBERO benchmark}
% Suites, levels, and the three metrics. Define CAR, TSR, ETS PRECISELY here,
% including the >1mm displacement threshold, and say it matches AEGIS so the
% comparison is fair. This paragraph is load-bearing for your whole evaluation.

\section{Obstacle geometry and the barrier}
% MVEE from the true mesh point cloud; the EE model; sphere decomposition.

\section{The CBF shield}
% The QP, the AEGIS-faithful every-step application, and the implementation
% choices that turned out to matter.

\section{Collision attribution}
% Raw displacement vs contact-graph attribution, and why the RAW flag is the
% metric reported (the attribution is a lower bound).

\section{Teachers}
\subsection{The shielded policy as its own teacher}
\subsection{A privileged scripted expert}
% Per-scene teacher profiles: grasp side by obstacle clearance, grasp target
% from real object geometry. This is the ablation, not the main line.

\section{Distillation}
% Flow-matching BC on shielded actions; the clean + shield-active filters;
% DAgger rounds and why they are capped at two.

% ============================================================================
\chapter{Experiments and Results}\label{ch:results}
% BUDGET ~9 pages. Organised by QUESTION, not by experiment: questions are stable, experiments
% accumulate under them. Each section follows the same template --- setup, table, how to read it,
% caveat --- so a new result is a fill-in rather than a restructure. Sections live in
% thesis/sections/ so adding one is a new file plus one \input line.

\section{Experimental setup}\label{sec:res:setup}
% STABLE — this section should not change as experiments are added.
% Cover: 24 scenes (3 suites x 2 levels x 4 tasks); inits 0--34 train / 35--49 held out;
% 5 rollouts per scene => n=120 per arm; horizon 300; 10 denoising steps; noise level 0
% (deterministic); Wilson 95\% intervals; TSR / CAR / ETS / CBF-activation definitions;
% and the step-matching rule (arms are compared at equal GRADIENT STEPS, not equal epochs,
% because episode length differs by teacher: 16{,}140 for r1 vs 15{,}988 for the planner student).

\section{Does the Shield Make \texorpdfstring{$\pi_{0.5}$}{pi 0.5} Safe?}\label{sec:shield_efficacy}

The first question is whether a control barrier function shield, applied as a runtime filter, makes
a pretrained VLA safe on this benchmark at all. Table~\ref{tab:shield_efficacy} and
Figure~\ref{fig:shield_efficacy} report the comparison.

\begin{table}[htbp]
\centering
\caption{Safety and task performance for base $\pi_{0.5}$ with and without the CBF runtime shield.
  Pooled over 24 scenes, $n = 120$ held-out rollouts per condition, Wilson 95\% intervals.}
\label{tab:shield_efficacy}
\begin{tabular}{lccc}
\toprule
\textbf{Policy} & \textbf{TSR (95\% CI)} & \textbf{Collision rate (95\% CI)} & \textbf{CAR} \\
\midrule
Base, no shield & 58.3\% [49.4, 66.8] & 82.5\% [74.7, 88.3] & 17.5\% \\
Base + shield   & 71.7\% [63.0, 79.0] & 13.3\% [\phantom{0}8.4, 20.6] & 86.7\% \\
\bottomrule
\end{tabular}
\end{table}

\begin{figure}[htbp]
  \centering
  \includegraphics[width=\linewidth]{figures/fig_shield_efficacy.pdf}
  \caption{Effect of the CBF runtime shield on base $\pi_{0.5}$, pooled over all 24 scenes.
    Error bars are Wilson 95\% intervals at $n = 120$ per condition.}
  \label{fig:shield_efficacy}
\end{figure}

Unshielded, $\pi_{0.5}$ displaces the obstacle in 82.5\% of episodes. With the shield active this
falls to 13.3\%, a reduction of 69 percentage points with disjoint confidence intervals.

Success also rises, from 58.3\% to 71.7\%. This is not incidental: a collision frequently derails
the episode, knocking the target object out of reach or displacing the goal, so preventing
collisions helps the policy complete the task. This point is worth holding onto, because in
Section~\ref{sec:internalise} the same shield applied to an already-avoidant policy has the
opposite effect on success, and the two observations are only consistent once the mechanism is
stated.

\paragraph{Caveat on privileged information.}
The barrier is constructed from ground-truth obstacle geometry taken directly from the simulator,
not from perception. A deployed system would have to estimate that geometry from sensors and would
inherit the resulting error. The numbers here should therefore be read as an upper bound on what
this class of shield achieves, with the perception gap removed by construction.

\paragraph{External validation.}
The evaluation pipeline can be checked against the figures reported for $\pi_{0.5}$ by the work
that introduced \texttt{SafeLIBERO}. Table~\ref{tab:external_validation} places the two side by
side. The unshielded rows agree to within a fraction of a point on both metrics, which supports the
view that this pipeline recovers the published baseline rather than producing an artefact of the
present implementation --- a check worth making explicitly, since every subsequent result in this
chapter is measured relative to that baseline.

The shielded rows diverge, and the divergence is worth stating rather than smoothing: the shield
used here is roughly fifteen points safer and three points less successful. Two deliberate
differences in the constraint account for it. The published formulation, by its authors' own
statement, ``solely constrains the end-effector'', and its limitations section notes that
unconstrained kinematic links may consequently collide; the barrier used here additionally
constrains links three through seven and the hand body (Section~\ref{sec:barrier}), and
Section~\ref{sec:scope} shows that arm links are exactly where residual collisions concentrate.
Separately, this barrier reads ground-truth geometry from the simulator, whereas the published
layer grounds its constraints through vision-language perception and fused depth, and attributes
its own residual collisions primarily to that upstream pipeline rather than to the controller.
Both differences favour the shield used here, so these figures are an upper bound for this class of
filter rather than a like-for-like improvement over that method.

% Table 4.1 — External validation against the benchmark's originating work (AEGIS/VLSA).
% Requires: \usepackage{booktabs}  \usepackage{threeparttable}
%
% NUMBERS AND THEIR PROVENANCE:
%   This work      — figures/results_shielded_distillation.md, eval 2026-08-07, n = 120 held-out.
%   AEGIS          — their Table I, averaged over the three suites used here (spatial, goal,
%                    object). Their own averages additionally include a Long suite, so these are
%                    recomputed three-suite figures, NOT their printed averages. Verified against
%                    the PDF 2026-08-18; see papers/FINDINGS.md.
%   Do NOT cite 87.5% for their shielded CAR. That figure came from a docstring in their released
%   code, not the paper, and is wrong. The paper's three-suite figure is 71.9%.

\begin{table}[htbp]
\centering
\begin{threeparttable}
\caption[External validation against the published baseline]{%
  External validation. The unshielded baseline measured here reproduces the published
  $\pi_{0.5}$ figures to within a fraction of a point on both metrics, indicating that the
  evaluation pipeline recovers the published result rather than an artefact of this
  implementation. The shielded rows are \emph{not} like-for-like and are discussed below.}
\label{tab:external_validation}
\begin{tabular}{llcc}
\toprule
\textbf{Condition} & \textbf{Source} & \textbf{CAR} & \textbf{TSR} \\
\midrule
\multirow{2}{*}{Base $\pi_{0.5}$, no shield}
  & Published\tnote{a}      & 17.3\%          & 58.9\%          \\
  & This work\tnote{b}      & \textbf{17.5\%} & \textbf{58.3\%} \\
\addlinespace[2pt]
\multicolumn{2}{l}{\emph{difference}} & $0.2$ pp & $0.6$ pp \\
\midrule
\multirow{2}{*}{With runtime CBF shield}
  & Published\tnote{a}      & 71.9\%          & 74.6\%          \\
  & This work\tnote{b}      & 86.7\%          & 71.7\%          \\
\addlinespace[2pt]
\multicolumn{2}{l}{\emph{difference}\tnote{c}} & $+14.8$ pp & $-2.9$ pp \\
\bottomrule
\end{tabular}
\begin{tablenotes}[flushleft]\footnotesize
\item[a] Averaged over the three suites used in this work (spatial, goal, object) from the
  originating work's reported per-suite results in the full action space. Their printed averages
  additionally include a Long suite and are therefore not directly quoted here.
\item[b] $n = 120$ held-out rollouts, pooled over 24 scenes. CAR is the complement of the collision
  rate defined in Section~\ref{sec:metrics}; note that the originating work defines collision
  avoidance as the proportion of strictly collision-free episodes without stating a numeric
  tolerance, so the 1\,mm threshold used here is this work's operationalisation of that definition.
\item[c] The shielded rows are not a like-for-like comparison and the gap is expected. Two
  differences in the constraint account for it, both deliberate: the published formulation
  constrains only the end-effector, by its authors' own statement, whereas the barrier used here
  additionally constrains links three to seven and the hand (Section~\ref{sec:barrier}); and this
  barrier reads ground-truth geometry from the simulator, whereas theirs is grounded through
  vision-language perception and fused depth, to which they attribute their own residual
  collisions. Both differences favour the shield used here, so these figures are an upper bound
  for this class of filter rather than an improvement over that method.
\end{tablenotes}
\end{threeparttable}
\end{table}


% =====================================================================================
%  Section 4.3 — Can the shield's safety be internalised?
%  Self-contained. Nothing to \input.
%
%  PREAMBLE (all already in main.tex): graphicx, booktabs, multirow, threeparttable
%  Additionally requires \usepackage{siunitx} ONLY if you switch to S columns; not used here.
%
%  ASSET: figures/fig_internalisation.pdf
%    Regenerate:  PYTHONPATH=. python -m experiments.make_fig_internalisation
%
%  STILL TO MAKE: Figure 4.3, the qualitative filmstrip. See the commented block at the end of
%  this file -- it cannot be built until the source episodes are identified.
%
%  NUMBER PROVENANCE: figures/results_shielded_distillation.md, eval 2026-08-07,
%  held-out inits 35-39, n = 120 per condition, 24 scenes.
% =====================================================================================

\section{Can the Shield's Safety Be Internalised?}\label{sec:res:internalise}

A runtime filter must run forever. It occupies inference budget on every control step, requires
obstacle geometry at deployment, and is a component that can fail. The central question of this
work is whether the behaviour it induces can instead be absorbed into the policy, so that the
policy is safe with the filter switched off.

Demonstrations were collected by rolling out the shielded policy and retaining episodes that
succeeded, displaced nothing, and in which the shield actually intervened at least once. That last
criterion matters: a clean episode in which the shield never fired is ordinary base-policy
behaviour and teaches nothing about safety. The retained set was 186 demonstrations from 480
rollouts, a 39\% yield. These were behaviour-cloned into the policy's action head, and the
procedure was repeated for a second round.

\begin{figure}[htbp]
  \centering
  \includegraphics[width=\linewidth]{figures/fig_internalisation.pdf}
  \caption[Internalisation of the shield's safety]{%
    The central result. The self-distilled policy, running with \emph{no shield at inference}
    (solid), is both safer and more successful than the base policy it was distilled from, and
    recovers most of the safety of the shielded baseline (hatched). Error bars are Wilson 95\%
    intervals at $n = 120$ held-out rollouts per condition.}
  \label{fig:internalisation}
\end{figure}

\begin{table}[htbp]
\centering
\begin{threeparttable}
\caption[Effect of self-distillation on safety and task performance]{%
  Effect of self-distillation, with and without the shield at inference. Rows in bold are the
  shield-free distilled policies, which are the claim of this work. $n = 120$ held-out rollouts
  per row, pooled over 24 scenes; intervals are Wilson 95\%.}
\label{tab:internalisation}
\begin{tabular}{lcccrr}
\toprule
\textbf{Policy} & \textbf{TSR (95\% CI)} & \textbf{Collision (95\% CI)} & \textbf{CAR}
  & $\lvert r_{\mathrm{cbf}}\rvert$\tnote{a} & \textbf{ETS}\tnote{b} \\
\midrule
Base, no shield        & 58.3\% [49.4, 66.8] & 82.5\% [74.7, 88.3] & 17.5\% & ---   & 138.6 \\
Base + shield          & 71.7\% [63.0, 79.0] & 13.3\% [\phantom{0}8.4, 20.6] & 86.7\% & 0.533 & 160.6 \\
\addlinespace[3pt]
\textbf{Round 1, no shield} & \textbf{82.5\% [74.7, 88.3]} & \textbf{19.2\% [13.1, 27.1]}
  & \textbf{80.8\%} & --- & 157.3 \\
Round 1 + shield       & 70.8\% [62.2, 78.2] & \phantom{0}8.3\% [\phantom{0}4.6, 14.7] & 91.7\% & 0.265 & 170.0 \\
\addlinespace[3pt]
\textbf{Round 2, no shield} & \textbf{80.8\% [72.9, 86.9]} & \textbf{17.5\% [11.7, 25.3]}
  & \textbf{82.5\%} & --- & 158.8 \\
Round 2 + shield       & 72.5\% [63.9, 79.7] & \phantom{0}9.2\% [\phantom{0}5.2, 15.7] & 90.8\% & 0.273 & 176.7 \\
\bottomrule
\end{tabular}
\begin{tablenotes}[flushleft]\footnotesize
\item[a] Mean absolute CBF reward term, a proxy for how much the shield had to correct. Necessarily
  zero when no shield is running, shown as ---.
\item[b] Execution time to success, in control steps, over successful rollouts only
  (Section~\ref{sec:metrics}). Conditional on success and must be read alongside TSR.
\end{tablenotes}
\end{threeparttable}
\end{table}

The headline result is the third row. With no shield running at inference, the distilled policy
displaces the obstacle in 19.2\% of episodes, against 82.5\% for the policy it was distilled
from~--- and its success rate rises from 58.3\% to 82.5\%. Both intervals are disjoint from the
base policy's. The distilled policy is simultaneously safer and more capable than its own teacher,
without the teacher's filter.

It is worth noting how close this comes to the shielded baseline: CAR 80.8\% shield-free against
86.7\% with the shield active. Most of the shield's safety benefit survives its removal.

\paragraph{One round is enough.}
Rounds one and two overlap on both metrics, so the second round buys nothing measurable. It also
does not erode performance, which is worth stating explicitly: an earlier DAgger-style experiment
in this project degraded across rounds (Section~\ref{sec:res:source}), and the absence of that
degradation here is informative rather than merely convenient.

\paragraph{Stacking the shield costs success.}
Applying the shield to the distilled policy improves safety further (CAR 80.8\% to 91.7\%) but
reduces success (82.5\% to 70.8\%). This is the mirror image of Section~\ref{sec:shield_efficacy}
and has the same explanation. The shield operates by projecting the policy's action onto the safe
set. When the policy frequently proposes unsafe actions, that projection is a net benefit. When the
policy already avoids the obstacle, the projection instead perturbs actions that were already both
safe and competent, displacing them from the distribution the policy was trained to produce~---
safe, but no longer a well-formed grasp. The claim of this work is therefore about shield-free
operation, not about stacking the two.

\paragraph{The cost of internalised caution.}
ETS rises from 138.6 to 157.3 steps, roughly 14\% slower. The distilled policy takes a more
conservative route. This is a real cost and is reported rather than omitted.

\paragraph{Characterising the training data.}
Because the demonstrations are the mechanism, their content is worth quantifying. Across the 186
retained demonstrations the shield was active on a median of 32\% of control steps, with a median
of 51 activations per episode, a minimum of nine, and no demonstration below five. The mean
correction norm was 0.33 against actions bounded at $\pm 1$ per axis. Roughly one imitated action
in three was therefore modified by the barrier, and modified substantially: the imitated action
distribution differs materially from the base policy's own.

% =====================================================================================
%  FIGURE 4.3 — qualitative filmstrip.  NOT YET BUILDABLE.
%
%  Intent: same scene and same held-out initial state, base pi0.5 (collides) above and the
%  distilled policy (avoids) below, 4-5 frames each, contact frame boxed or arrowed. Every other
%  figure in this chapter reports a rate; none shows the behaviour, so this is the largest
%  qualitative gain available to the chapter.
%
%  BLOCKED ON: videos/ contains task3_ep0.mp4, task3_ep1.mp4, task3_ep3.mp4 with no record of
%  which policy produced them, which scene they are, or whether any contains a collision. Pairing
%  them by guesswork would risk captioning a distilled rollout as the base policy.
%
%  TO UNBLOCK, record for two rollouts of the SAME scene and SAME held-out init:
%    (i) base policy, an episode that collides;  (ii) round-1 distilled, an episode that does not.
%  Then:  PYTHONPATH=. python -m experiments.make_fig_filmstrip \
%             --base <base.mp4> --distilled <distilled.mp4> --contact-frame <N>
%  (that script does not exist yet either -- write it once the inputs are known.)
%
%  \begin{figure}[htbp]
%    \centering
%    \includegraphics[width=\linewidth]{figures/fig_filmstrip.pdf}
%    \caption{...}
%    \label{fig:filmstrip}
%  \end{figure}
% =====================================================================================

\section{Does the source of the demonstrations matter?}\label{sec:res:source}
% THE BOUNDARY RESULT. Same benchmark, same student, same eval code, same held-out inits,
% matched gradient budget (15{,}988 vs 16{,}140 steps).
%
% base+noCBF       TSR 58.3 [49.4, 66.8]   collision 82.5 [74.7, 88.3]
% planner_A+noCBF  TSR 54.2 [45.3, 62.8]   collision 80.0 [72.0, 86.2]   <- flat against base
% r1+noCBF         TSR 82.5 [74.7, 88.3]   collision 19.2 [13.1, 27.1]
%
% Claim: distilling safety from a FOREIGN expert fails; distilling it from the policy's own
% shielded rollouts succeeds. State the caveat yourself: 6 of 24 scenes had zero planner demos,
% so ~30 of 120 rollouts measure teacher COVERAGE rather than distillation --- report the
% 18-scene slice alongside the 24-scene headline. It does not rescue the result.
%
% Mechanism: NOT observation aliasing (measured and refuted, see Sec.~\ref{sec:res:mechanism}).
% Surviving candidates: distribution shift, and action-distribution mismatch (the planner emits
% saturated P-control/QP deltas against \pizero{}'s flow-matched chunks).
%
% SLOT: randomised-teacher ablation (shielded teacher + randomised layouts) goes here when run.

% =====================================================================================
%  Section 4.5 — What can the shield not reach?
%  Self-contained. Nothing to \input.
%
%  PREAMBLE (already in main.tex): graphicx, booktabs, multirow, threeparttable
%  ASSET: figures/fig_culprits.pdf
%    Regenerate:  PYTHONPATH=. python -m experiments.make_fig_culprits
%
%  TWO THINGS NOT TO EDIT AWAY:
%   (1) The culprit categories OVERLAP -- base is 71+17+36 = 124 attributions over 99 collided
%       episodes. Any phrasing implying the columns partition the collisions is wrong, and the
%       figure is grouped rather than stacked for the same reason.
%   (2) The arm-link improvement (17 -> 8) has OVERLAPPING Wilson intervals and is stated as
%       suggestive. The gripper result (71 -> 0 across 360 shielded episodes) is not in doubt.
%       Do not let the two be read at the same confidence.
% =====================================================================================

\section{What Can the Shield Not Reach?}\label{sec:res:scope}

The shield reduces collisions to 13.3\% but not to zero. This section identifies what the residual
consists of, using per-body collision attribution recovered from the evaluation logs. Contacts
attributed to other scene objects are excluded throughout: that category fires on every episode of
every condition, including collision-free ones, because an obstacle resting on its supporting
surface registers a permanent contact.

\begin{table}[htbp]
\centering
\begin{threeparttable}
\caption[Residual collisions by culprit body]{%
  Per-body attribution of collisions. Counts are episodes out of $n = 120$ in which the named body
  was implicated. Categories are not mutually exclusive, so a row need not sum to its collided
  count.\tnote{a}}
\label{tab:culprits}
\begin{tabular}{lrrrrr}
\toprule
\textbf{Policy} & \textbf{Episodes} & \textbf{Collided} & \textbf{Gripper} & \textbf{Arm link}
  & \textbf{Held object} \\
\midrule
Base, no shield        & 120 & 99 & 71          & 17 & 36 \\
Base + shield          & 120 & 16 & \textbf{0}  & 15 & \phantom{0}0 \\
\addlinespace[2pt]
Round 1, no shield     & 120 & 23 & \phantom{0}7 & 13 & \phantom{0}9 \\
Round 1 + shield       & 120 & 10 & \textbf{0}  & \phantom{0}8 & \phantom{0}1 \\
\addlinespace[2pt]
Round 2, no shield     & 120 & 21 & \phantom{0}8 & \phantom{0}8 & \phantom{0}7 \\
Round 2 + shield       & 120 & 11 & \textbf{0}  & \phantom{0}9 & \phantom{0}0 \\
\bottomrule
\end{tabular}
\begin{tablenotes}[flushleft]\footnotesize
\item[a] For example the base row records 124 attributions across 99 collided episodes, because a
  single episode may implicate more than one body.
\end{tablenotes}
\end{threeparttable}
\end{table}

\begin{figure}[htbp]
  \centering
  \includegraphics[width=\linewidth]{figures/fig_culprits.pdf}
  \caption[Collision attribution by body across conditions]{%
    What the shield reaches. The end-effector channel is eliminated wherever the shield is running
    (shaded columns): zero gripper collisions across all 360 shielded episodes. Arm-link contacts
    survive it almost unchanged. Bars are grouped rather than stacked because the categories
    overlap, and the dotted line marks the number of episodes with any collision at all.}
  \label{fig:culprits}
\end{figure}

\paragraph{The shield eliminates the end-effector channel completely.}
Gripper collisions fall from 71 to zero, and across all 360 shielded episodes there is not a single
one. Of the 16 residual collisions in the shielded baseline, 15 involve an arm link.

\paragraph{The residual is a fidelity limit, not a scope limit.}
Arm links are not absent from the barrier: links three to seven and the hand are constrained
(Section~\ref{sec:barrier}). But they are represented far more coarsely than the end-effector ---
three point samples per link carrying a single sphere radius, against three fitted spheres for the
hand and fingers. More importantly, an arm link's velocity is not the quantity the quadratic
program optimises. The program solves for an end-effector velocity and infers link motion through a
damped Jacobian pseudo-inverse, so the constraint is enforced against a \emph{predicted} link
velocity that diverges from the realised one wherever the inverse is ill-conditioned or the
controller's null-space motion differs from the minimum-norm solution.

The pattern in Table~\ref{tab:culprits} is therefore the expected one: collisions are eliminated
exactly where the barrier's model of the robot is precise and its velocity is the controlled
variable, and they persist where the model is coarse and the velocity is estimated. The corollary
for reading Section~\ref{sec:res:internalise} still holds --- the distilled policy's shield-free
figure is within a few points of what its teacher actually achieved, not a degraded approximation
of a perfect one --- but the teacher's own ceiling is set by approximation error in the constraint,
not by an absence of authority.

\paragraph{A second transfer channel.}
Arm-link collisions fall from 17 to 8 between the base and twice-distilled policies, while the
shielded baseline --- with the arm-link constraints active --- still shows 15. The improvement is
therefore not inherited from the barrier's own arm-link performance, which is worse than the
student's. It must come instead from the selection criterion: demonstrations were retained only if
nothing at all was displaced, arm links included, and the student learned that. Safety appears to
transfer through two distinguishable channels --- imitation of the filter's corrections where the
filter has authority, and outcome selection where it does not.

This second channel is stated with less confidence than the first. As proportions, 17/120 against
8/120 is 14.2\% [9.0, 21.5] versus 6.7\% [3.4, 12.6]; those Wilson intervals overlap, so the
arm-link trend is suggestive rather than established. The gripper result is not in doubt.

\paragraph{Independent corroboration.}
A separate experiment measured the end-effector's closest approach to the obstacle during level-II
episodes that were scored as collisions. The minimum was 0.231\,m against a barrier radius of
0.15\,m: the end-effector never came close, yet the episode collided. Together with the attribution
table and the inert guided-sampling result of Section~\ref{sec:res:channels}, three independent
measurements point at the same conclusion.

% =====================================================================================
%  Section 4.6 — What explains the improvement?
%  Self-contained. Nothing to \input.
%
%  PREAMBLE (already in main.tex): graphicx, booktabs, multirow, threeparttable
%  ASSETS: figures/fig_matched_control.pdf, figures/fig_aliasing.pdf
%    Regenerate both:  PYTHONPATH=. python -m experiments.make_fig_mechanism
%
%  DO NOT SOFTEN THE ALIASING SUBSECTION. It reports a prediction this project made and then
%  refuted with its own data, and it records that the prediction was written down before the
%  planner figure was computed. That ordering is what makes it evidence rather than a
%  rationalisation, and it is worth more to the examiner than a result that merely worked.
% =====================================================================================

\section{What Explains the Improvement?}\label{sec:res:mechanism}

Section~\ref{sec:res:internalise} showed that distilling the shielded policy's own clean rollouts
makes it dramatically safer without the shield. It does not, on its own, establish \emph{why}. The
demonstrations were filtered on three properties at once --- the shield fired, the episode
succeeded, and nothing was displaced --- and two mechanisms are bundled inside that filter.

Under the first, \textbf{distillation}, the policy imitates the shield's avoidance corrections.
This is the claim of this work. Under the second, \textbf{selection}, the improvement comes from
training a policy on its own successful episodes, a known self-improvement effect requiring no
shield at all; successful episodes also happen to be the ones in which nothing was knocked over. If
the second mechanism accounts for the result, the shield is incidental and the contribution is a
rediscovery of filtered behaviour cloning.

Nothing in the six-condition evaluation of Section~\ref{sec:res:internalise} distinguishes them, so
a matched control was run.

\subsection*{The matched control}

Demonstrations were collected with the shield \textbf{switched off}, filtered by exactly the same
criteria --- succeeded and displaced nothing --- and used to train a student with identical
hyperparameters. The shielded pool was then subsampled to the same size, so the two conditions
differ in one respect only: whether the shield was active while the demonstrations were being
generated.

\begin{table}[htbp]
\centering
\caption{The matched control. Both distilled conditions use 85 demonstrations passing an identical
  filter; they differ only in whether the shield was active during collection. $n = 120$ held-out
  rollouts per row, Wilson 95\% intervals.}
\label{tab:matched_control}
\begin{tabular}{lccc}
\toprule
\textbf{Condition} & \textbf{Demos} & \textbf{TSR (95\% CI)} & \textbf{Collision (95\% CI)} \\
\midrule
Undistilled base & ---            & 58.3\% [49.4, 66.8] & 82.5\% [74.7, 88.3] \\
Control, shield \textbf{off}  & 85 & 50.0\% [41.2, 58.8] & 84.2\% [76.6, 89.6] \\
Matched, shield \textbf{on}   & 85 & \textbf{75.8\% [67.4, 82.6]} & \textbf{26.7\% [19.6, 35.2]} \\
\bottomrule
\end{tabular}
\end{table}

\begin{figure}[htbp]
  \centering
  \includegraphics[width=\linewidth]{figures/fig_matched_control.pdf}
  \caption[The matched control: selection versus correction]{%
    Separating selection from correction. Training on the policy's own successful, collision-free
    episodes (centre) leaves it indistinguishable from the undistilled baseline. Applying the same
    filter to demonstrations collected while the shield was active (right) produces the effect.
    Since the two distilled conditions use the same number of demonstrations, the same filter and
    the same hyperparameters, the difference is attributable to the shield.}
  \label{fig:matched_control}
\end{figure}

\paragraph{Success-filtering alone achieves nothing.}
The control sits at 84.2\% collision, statistically indistinguishable from the undistilled base
policy at 82.5\%. Training a policy on its own successful, collision-free episodes did not make it
safer on this benchmark.

\paragraph{The shield's corrections are what transfers.}
At the same demonstration count and under the same filter, the shielded condition reaches 26.7\%,
with a confidence interval disjoint from the control's. The difference between the two conditions
is attributable to the shield and to nothing else in the pipeline. This is the experiment that
answers the most natural objection to Section~\ref{sec:res:internalise}, and it answers it directly
rather than by argument.

\paragraph{Two further pieces of evidence.}
The conclusion is corroborated by two measurements taken for other purposes. First, the CBF
activation proxy falls from 0.533 to 0.265 between the shielded base and the shielded distilled
policy: after distillation the shield finds materially less to correct. Success-filtering alone
does not predict this --- a policy that merely completes tasks more often would not require
\emph{less} safety intervention. Second, the demonstrations are pervasively shield-shaped, with the
barrier active on a median 32\% of control steps and a mean correction norm of 0.33 against actions
bounded at $\pm 1$, so the imitated action distribution differs substantially from the base
policy's own.

\paragraph{A nuance worth preserving.}
The control is not uniformly inert. Its per-body attribution shows arm-link collisions falling from
17 to 6 while gripper collisions remained high at 44, which is why the total did not move. Outcome
selection therefore \emph{does} transfer avoidance in the channel the barrier cannot express ---
exactly the second transfer channel identified in Section~\ref{sec:res:scope} --- but it is not
sufficient on its own. The two mechanisms are complementary rather than competing, and only one of
them addresses the end-effector channel that dominates the collision count.

\paragraph{The cost of collecting without a shield.}
One further number deserves reporting. The unshielded collection retained 85 usable demonstrations
from 480 rollouts, a yield of 18\%, against 39\% for the shielded collection. Safe demonstrations
are roughly twice as expensive to obtain without a shield. This is a second, independent sense in
which the shield earns its place: beyond supplying the corrections that transfer, it makes the
collection of a safe demonstration set tractable in the first place.

\subsection*{Two alternative explanations tested and rejected}

Section~\ref{sec:res:source} established that distilling from a foreign expert fails while
self-distillation succeeds. Two explanations for that boundary were tested directly, and both were
rejected.

\paragraph{Observation aliasing.}
Earlier analysis in this project attributed the failure to the policy's observation space:
$\pi_{0.5}$ sees an image and eight proprioceptive dimensions, with no joint angles, so a
correction depending on the full arm configuration is not a function of anything the policy
observes, and behaviour cloning would average contradictory targets. This was asserted but never
measured, and it is measurable from the demonstration files alone. For each dataset, samples whose
observations are near neighbours were found and the disagreement between their target actions
measured, normalised by the disagreement between random pairs. Comparison is made at matched
observation \emph{radius} rather than matched neighbour count, because a nearest-neighbour ratio
conflates aliasing with sampling density and the planner's trajectories are far denser.

\begin{figure}[htbp]
  \centering
  \includegraphics[width=0.72\linewidth]{figures/fig_aliasing.pdf}
  \caption[Observation aliasing does not explain the transfer failure]{%
    A prediction of this project, refuted by its own data. If the planner's corrections were
    poorly determined by what the policy observes, its curve would lie \emph{above} the shielded
    one. It lies consistently below, by a factor of roughly two at every radius: the teacher that
    fails to distil has the cleaner observation-to-action mapping.}
  \label{fig:aliasing}
\end{figure}

\begin{table}[htbp]
\centering
\caption{Action disagreement among observation-neighbours, normalised by random-pair disagreement.
  Lower values mean the action is better determined by the observation.}
\label{tab:aliasing}
\begin{tabular}{lcc}
\toprule
\textbf{Radius (z-scored)} & \textbf{Shielded (distils)} & \textbf{Planner (does not)} \\
\midrule
0.10 & 0.103 & 0.051 \\
0.25 & 0.174 & 0.093 \\
0.50 & 0.237 & 0.118 \\
1.00 & 0.376 & 0.174 \\
\bottomrule
\end{tabular}
\end{table}

The ordering is the reverse of the prediction. The planner's demonstrations are roughly twice as
\emph{well} determined by the observation at every radius: the teacher that fails to distil has the
cleaner observation-to-action mapping, and the teacher that succeeds is the more aliased of the
two. In hindsight the direction is unsurprising --- a scripted planner is close to a deterministic
function of state --- but it means aliasing is not the binding constraint, and the prior
explanation does not survive contact with the data. The prediction was recorded before the planner
figure was computed. The surviving candidates are distribution shift and action-distribution
mismatch, and Section~\ref{sec:res:source} takes up the first of them.

\paragraph{Generative uncertainty as a collision predictor.}
$\pi_{0.5}$ produces each action by integrating a flow ordinary differential equation, and every
rollout stores that denoising chain. If the policy were less certain in states where safe and
unsafe behaviours are both plausible, the chain geometry would carry a risk signal, and a runtime
monitor could use it with no barrier and no obstacle geometry. Scored as the area under the ROC
curve for predicting whether an episode collided, the result is negative: 0.461 for the base
policy, which is chance. The distilled policy shows a consistent inverted association across all
three chain statistics (AUC 0.350, equivalently 0.65 when flipped, so \emph{lower} uncertainty is
associated with collision), but on 23 collision episodes that is roughly two standard errors from
chance and far short of deployable. The suggestive reading, if it is real, is that the distilled
policy's failures are \emph{confident} ones, which would make uncertainty-based monitoring
structurally unsuited to this failure mode.

% =====================================================================================
%  Section 4.7 — Are there other channels for safety?
%  Self-contained. Nothing to \input.
%
%  PREAMBLE (already in main.tex): graphicx, booktabs, multirow, threeparttable
%  ASSETS: figures/fig_rl_bracket.pdf, figures/fig_bon_correlation.pdf
%    Regenerate both:  PYTHONPATH=. python -m experiments.make_fig_channels
%
%  THE SUMMARY TABLE IS THE POINT OF THIS SECTION. Four negatives spread over several pages
%  read as a list of things that did not work; the same four with their tested scope stated
%  read as a bounded result. If space forces a cut, cut prose and keep Table 4.7.
%
%  EVERY NEGATIVE HERE IS SCOPED. Do not let any of them be restated as a general claim:
%  the RL result is about the reward design tested, best-of-N is about this policy and
%  benchmark, guided sampling is about an end-effector barrier, and the language result is
%  about the phrasing tested.
% =====================================================================================

\section{Are There Other Channels for Safety?}\label{sec:res:channels}

Four further mechanisms were tested and none replaced the shielded-imitation route. They are
reported together because a characterised negative bounds the claim of this work in a way a
positive result cannot, and because each fails for a different and identifiable reason.
Table~\ref{tab:negatives} summarises them; the subsections give the evidence.

\begin{table}[htbp]
\centering
\begin{threeparttable}
\caption[Mechanisms tested and rejected]{%
  Mechanisms tested and rejected, with the scope at which each negative is claimed. None is a
  general claim about its method class.}
\label{tab:negatives}
\footnotesize
\begin{tabular}{p{0.17\linewidth}p{0.24\linewidth}p{0.22\linewidth}p{0.27\linewidth}}
\toprule
\textbf{Mechanism} & \textbf{What was tested} & \textbf{Outcome} & \textbf{Scope of the negative} \\
\midrule
Scalar-reward RL
  & Flow-SDE GRPO on the action head, 4 configurations $\times$ 6 rounds, one scene
  & Either collapses to inaction or learns nothing; no setting does both
  & The reward design and budget tested. Not safe RL as a class; constrained-MDP and
    world-model methods supply richer signals \\
\addlinespace[3pt]
Language-expressed constraint
  & One safety clause appended to the instruction, $n=60$ per condition, level II
  & No safety change; 23\% slower, 11.7 points less successful
  & The single phrasing tested: generic, unnamed referent, manner adverbial \\
\addlinespace[3pt]
CBF-guided sampling
  & Barrier gradient injected into the flow integration, validated on a synthetic barrier
  & Never activates on real scenes; closest approach 0.231\,m against a 0.15\,m barrier
  & Guidance built on an \emph{end-effector} barrier. Untested with the full arm-link
    constraint set \\
\addlinespace[3pt]
Best-of-$N$ selection
  & $K \in \{4, 8\}$ candidates re-simulated per query, exact state rewind
  & No safe candidate on $\sim$60\% of queries; doubling $K$ buys 8.8 points
  & This policy and benchmark. The measurement it yields is reported as the result \\
\bottomrule
\end{tabular}
\begin{tablenotes}[flushleft]\footnotesize
\item Two further \emph{explanations}, as opposed to mechanisms, were also tested and rejected in
  Section~\ref{sec:res:mechanism}: observation aliasing and generative uncertainty.
\end{tablenotes}
\end{threeparttable}
\end{table}

\subsection*{Scalar-reward reinforcement learning}

Reinforcement learning was tried first, and its failure is what motivated the imitation approach
that the rest of this chapter evaluates. It is reported at length because the failure is
structural, and because four runs are needed to establish that.

The setup was flow-SDE GRPO on $\pi_{0.5}$'s flow-matching action head: the deterministic sampling
ODE is converted to an equivalent SDE so that log-probabilities exist and a policy gradient can be
taken. A LoRA adapter on the action head received gradients with the backbone frozen, matching the
trainable surface used for distillation in Section~\ref{sec:res:internalise}, so the comparison
between the two learning signals is not confounded by capacity. Rollouts used the CPS sampler at
noise 0.7 with ten denoising steps and clipping 0.2, on one level-II scene, with $K = 8$ rollouts
per group over six rounds. The reward combined success, direct collision, CBF activation rate and
progress terms.

\begin{figure}[htbp]
  \centering
  \includegraphics[width=\linewidth]{figures/fig_rl_bracket.pdf}
  \caption[Scalar-reward RL bracketed from both sides]{%
    The two informative configurations. \textbf{Left:} admitting real collisions to the reward
    moves the policy, and it collapses --- success reaches zero \emph{in the shielded condition
    too}, which the shield alone cannot cause, so the policy has found the trivial zero-collision
    optimum of doing nothing. \textbf{Right:} a lower learning rate with an annealed shield fixes
    the collapse and learns nothing, with unshielded collision pinned at 1.0 for six rounds.
    Neither panel alone is an argument; together they are the two sides of one knob.}
  \label{fig:rl_bracket}
\end{figure}

\paragraph{Too weak.}
With the shield active on every rollout, collisions never enter the reward directly --- the shield
prevents them --- so the only safety signal is the CBF activation penalty. Raising the learning
rate to $5 \times 10^{-5}$ and the activation weight to 1.5 left the policy flat.

\paragraph{Too strong.}
Running half of each group's rollouts unshielded puts real collisions into the reward, widening the
within-group spread to roughly $\Delta = 1.0$. The safety metric improved immediately and the task
collapsed with it. Two causes compound: the adapter diverged, since the mixed reward is far higher
variance than the flat one at the same learning rate; and the collision penalty punishes progress,
because in these scenes the obstacle lies between the gripper and the goal, so nearly all forward
motion collides early in training.

\paragraph{What the bracket establishes.}
The only configuration that moved the policy's behaviour collapsed it; every configuration stable
enough to preserve the task learned no avoidance. That is not a gap between two tuned settings but
the two sides of a single one. The learning signal and the destabilising signal are the same knob,
for two compounding reasons. A single scalar episode reward gives no per-step spatial credit, so
the policy cannot separate ``detour around the obstacle'' from ``stop moving toward the goal'' when
the obstacle lies on the path to the goal. And in a flow-matching policy, exploration \emph{is}
sampling noise, which is also what degrades the actions being evaluated.

This is precisely what the shield supplies and the reward does not. The barrier already computes
the correct action at every step it intervenes, in the same space the policy emits, so the credit
assignment the reward could not infer is simply given.

\subsection*{Safety expressed in language}

Every mechanism examined so far acts on the action. A vision-language-action model exposes a
channel none of them use, and one a classical controller does not possess at all --- the
instruction itself.

The experiment is minimal by design: the same checkpoint, seed, held-out initial states and
evaluation code for both conditions, differing only in a clause appended to the instruction. Twelve
level-II scenes at five initial states each give $n = 60$ per condition; level II was chosen
because collisions concentrate there.

\begin{quote}\small
\textbf{plain:} \texttt{pick up the black bowl between the plate and the ramekin and place it on
the plate}\\[3pt]
\textbf{safety:} \texttt{\dots{} and place it on the plate, being careful not to touch or knock
over any other objects}
\end{quote}

The suffix is reproduced verbatim because a negative about language is only as strong as the
language tested. The phrasing is generic: it names no object, gives no spatial reference, and
states the constraint as a manner adverbial rather than a prohibition. Each is a design choice that
could plausibly matter, and none was varied.

\begin{table}[htbp]
\centering
\caption{Effect of appending a safety clause to the instruction. $n = 60$ per condition, twelve
  level-II scenes, Wilson 95\% intervals.}
\label{tab:language}
\begin{tabular}{lccc}
\toprule
\textbf{Condition} & \textbf{TSR (95\% CI)} & \textbf{Collision (95\% CI)} & \textbf{ETS} \\
\midrule
Plain instruction       & 66.7\% [54.1, 77.3] & 70.0\% [57.5, 80.1] & 132.4 \\
Safety clause appended  & 55.0\% [42.5, 66.9] & 68.3\% [55.8, 78.7] & 162.4 \\
\bottomrule
\end{tabular}
\end{table}

Collisions fall by 1.7 percentage points, with almost coincident intervals. But the policy is not
ignoring the text: execution time rises 23\% and success falls 11.7 points, while the per-body
attribution shows no meaningful redistribution. The policy responds to the safety clause by acting
more slowly and completing fewer tasks, without touching anything less often.

The natural reading is that safety language elicits a general behavioural prior --- hesitancy,
reduced speed --- rather than a grounded avoidance behaviour. The model has learned that
safety-flavoured instructions correspond to cautious-looking motion, and that association is not
connected to the geometry of the scene. It produces the appearance of caution without its
substance. The distinction matters: it separates ``the model ignores safety language'' from ``the
model responds to safety language but does not ground it'', and only the second says anything about
what such a channel would need.

\subsection*{CBF-guided sampling}

The shield corrects an action after the policy has produced it. A flow-matching policy offers an
earlier point of intervention: the action is generated by integrating an ODE over several denoising
steps, so a barrier gradient can be injected \emph{into} that integration, steering generation
toward safe actions rather than projecting unsafe ones afterwards.

The implementation was validated before deployment, because two sign and scheduling conventions
produce plausible-looking but incorrect behaviour. On a synthetic barrier with a known closed-form
projection, $\lambda = 1$ reproduces the projection endpoint exactly.

On real scenes it never fired. Across level-II episodes the end-effector's closest approach to the
obstacle was 0.231\,m against a barrier radius of 0.15\,m --- the guidance term was identically
zero for the entire rollout --- while those same episodes collided. The reason is the one
Section~\ref{sec:res:scope} identifies: the barrier is a function of the end-effector, and the
contacts that occur are made by the arm links. Guidance inherits the scope of the barrier it steers
by, so steering generation cannot reach a failure the barrier cannot see. This is a scope-limited
negative, and it would be worth revisiting on a barrier carrying the full arm-link constraint set.

\subsection*{Best-of-$N$ selection, and what it measured instead}

If safety varied across the actions a policy might sample from a given state, it could be obtained
by rejection: draw $K$ candidate action chunks, simulate each, execute the safest. The sampler must
be stochastic for candidates to differ, so the matched control is $K = 1$ at the same noise level
rather than the deterministic policy used elsewhere in this chapter. Making the measurement
trustworthy required the simulator to be rewound exactly between candidates; snapshot and restore
were verified to machine precision before any result was recorded.

\begin{figure}[htbp]
  \centering
  \includegraphics[width=0.82\linewidth]{figures/fig_bon_correlation.pdf}
  \caption[Within-state safety correlation in a VLA policy]{%
    Why rejection sampling cannot help here. In every configuration the fraction of individually
    safe candidates coincides with the fraction of queries in which \emph{all} $K$ are safe, so a
    query is $0$-of-$K$ or all-of-$K$ safe and almost never in between. Independent sampling would
    make the all-$K$-safe rate two to three orders of magnitude smaller (right-hand bars).}
  \label{fig:bon}
\end{figure}

As a method this is negative: on roughly 60\% of queries no candidate was safe, and buying more
candidates barely moved it --- doubling $K$ from four to eight recovered 8.8 percentage points for
twice the compute, and raising the noise recovered 3.7.

The measurement underneath it is the more interesting result. The mean safe fraction and the
all-$K$-safe fraction coincide in all three configurations. Were the candidates independent draws,
the all-$K$-safe rate would be $p^{K}$: 1.1\% at $K = 4$ and 0.08\% at $K = 8$, against observed
values of 32\% and 41\%. The policy's samples from a given state are therefore near-perfectly
correlated in their safety outcome.

Safety in this policy is consequently \textbf{state-determined rather than sample-determined}, at
least on this benchmark. Once a state is reached the outcome is largely fixed, and no amount of
choosing among actions recovers it --- the leverage lies in not reaching the state. This also
explains a result that would otherwise look inconsistent: episode-level filtering worked
(Section~\ref{sec:res:mechanism}) while action-level selection cannot, because episode-level
selection chooses between trajectories that reached \emph{different states}, which is the axis
along which variation actually exists.

The correlation measurement is offered here as the contribution rather than the selection method.
It is, as far as this work is aware, the first measurement of candidate-level safety correlation in
a vision-language-action policy, and it is the sharpest available statement of the reachability
theme running through Sections~\ref{sec:res:scope}, \ref{sec:res:mechanism} and the discussion that
follows.

\section{Two defects in the benchmark}\label{sec:res:benchmark}
% A service contribution: both affect anyone else's SafeLIBERO level-II numbers.
%
%  1. NON-DETERMINISM. SafeLIBERO parks unused obstacles off-scene in a heap. Left collidable they
%     overflow MuJoCo's contact buffer (ncon = 5000), and on overflow MuJoCo drops contacts
%     order-dependently, so identical episodes scored differently in 6 of 26 re-runs (23\%).
%     Fixed by clearing contype/conaffinity on parked bodies.
%  2. PHANTOM COLLISIONS. Objects spawning unsupported fall under gravity with the arm stationary
%     --- a moka pot fell 102\,mm --- and the displacement is attributed to the robot. Fixed with
%     60 settling steps before measurement begins; residual 0.00\,mm.
%     Ordering matters: the settle must run AFTER parked contacts are disabled.


\chapter{Discussion}\label{ch:discussion}
% BUDGET ~3 pages. "Critical analysis of method and results" is an explicit
% marking criterion — this chapter is worth real marks. Do not skip the
% limitations; naming them yourself reads as strength, not weakness.

\section{Why imitation succeeds where scalar-reward RL fails}
\section{The cost of internalised caution}
% Success erodes across DAgger rounds. Explain the mechanism.
\section{Limitations}
% Simulation only; privileged geometry; OSC_POSE action-space limits on the
% wedged-object tasks; single seed where applicable.
\section{Future work}

% ============================================================================
\chapter{Conclusion}\label{ch:conclusion}
% BUDGET ~1 page. Short. Nothing new.

% ============================================================================
\bibliographystyle{plainnat}
\bibliography{references}

\appendix
\chapter{Additional results}
% Per-suite and per-task breakdowns that would clutter the main text.

\chapter{Implementation notes}
% Reproduction commands. NOTE: the module guidance requires crediting AI tools
% used for code or writing, both in the codebase README and here. Be specific
% about which parts.

\end{document}
